% Modernised reading — fonds Alexandre Grothendieck, shelfmark 4, pages 1-11.
%
% This is an INTERPRETATION, not a transcription. Everything the critical
% apparatus carried moves into footnotes. In any dispute of substance the
% transcription governs: whoever cites Grothendieck must cite batch-01.fr.tex.
%
% Pass: Opus 5.5 (claude-opus-5-5), 2026-10-03 — first pass, unchecked against
% the pages by a human. The folder was read whole; it holds one batch, pages
% 1-11, all transcribed.
%
% What the folder is. §8 of a longer typescript, « Une estimation
% (heuristique) du nombre des valeurs propres de frobenius dans la cohomologie
% évanescente qui sont des unités p-adiques » (pages 2-11), which he revised by
% hand, and a manuscript leaf « 8.4.4 » (page 1) written for that revision.
% The reading follows the typescript's order and reads page 1 where the
% revision puts it, after Lemme 8.4.3.
%
% The revision changes the logic, and the reading is organised around that.
% The typed text computes the torsion terms of the snake sequence through
% Poincaré duality (8.2, 8.3.2, 8.5.3) and so takes F-semisimple parts of the
% duals; the revision (page 1, margins of pages 4-6) computes them directly:
% (Coker u_2)_ss = 0 and (Ker u_2)_ss = (pTors^n(X))_ss, by the nilpotence of
% F on p-torsion above the dimension. Hence phi is always surjective, and the
% typed claims and examples of non-surjectivity (pages 9-10) are struck. The
% reading takes the revised logic as the text, gives the typed one as what was
% abandoned, and says (ours) why: duality exchanges F with its transpose, so an
% F-semisimple part is dual to a semisimple part for the transpose, not for F.
%
% Notation fixed for the whole document. k perfect of char. p > 0, W = W(k),
% sigma its Frobenius; X smooth projective connected of dimension n in P^r_k;
% H^i(X) the crystalline cohomology over W; Tors^i(X) its torsion,
% H^i_lib(X) = H^i(X)/Tors^i(X), {}_pTors^i(X) = Tors^i(X)[p] =
% Tor_1^W(H^i(X), k). Y a smooth section of X by a hypersurface of degree d.
% For a finite-dim k-space with a p-linear F: V_ss (subscript) its
% F-semisimple part; superscript ss for the transpose operator V. O_X is
% written \mathcal{O}_X (underlined on the typescript).
%
% Substantive departures from the pages, each footnoted where it occurs:
% — page 2, (8.1.3): « Ker u_0 » is Ker(u_0 ⊗ k); written so.
% — page 2 margin, second sequence: the last term's exponent, read i-n+1
%   (uncertain), must be i+1, as (**) of page 1 shows at i = n.
% — page 4-6: w is the hand's w of (8.5.1), Ker v -> (pTors^n)_ss, described
%   explicitly; the typed w of (8.5.3), struck, is not used.
% — page 5, Théorème 8.5: Poincaré duality is listed among the hypotheses but
%   the revised argument does not use it.
% — page 8, 8.5.9: « H^n(X, O_X)_ss » (exponent rewritten, uncertain) — what
%   serves is (pTors^n(X))_ss, a quotient of H^{n-1}(X, O_X)_ss.
% — page 9, 8.5.10: φ injective iff w surjective holds unconditionally; the
%   reserve Ker v = 0 belongs to the criterion by pTors^n.
% — page 10, (8.5.11.1): equality is equivalent to φ bijective only together
%   with (Tors^{n-1} ⊗ k)_ss = 0 (or if the torsion is not removed from the
%   right-hand side); stated so.
% — page 10, Serre's example: the page is cut off at « ≠ Z/ »; the reading
%   gives H^1(X̄, Z/p) ≅ Z/p, which is what the example has.
%
% Modern names given and footnoted as ours: théorème de comparaison de
% Berthelot-Ogus, décomposition de Fitting, correspondance de Lang/Katz
% (modules de Frobenius étales), opérateur de Cartier, ordinarité, théorème de
% Katz-Messing, formule des traces modulo p (« de Woods Hole »), Mazur
% (Newton au-dessus de Hodge), Illusie (ordinarité générique), comparaison
% entière de Bhatt-Morrow-Scholze.
%
% What the folder announces and never establishes: the injectivity of v in
% general (8.5.7, « je pense »); the non-vanishing of the generic
% E(Y, O_Y)_ss for d large (8.5.8, referred to n° 7, not in the folder); the
% conjecture of n° 6 on the level of evanescent cohomology; the monodromy
% question of 8.6 and the congruence for card Y(k'). Weak Lefschetz with
% torsion in crystalline cohomology is admitted throughout.
\documentclass[11pt,a4paper]{article}
\input{../preamble/grothendieck}

\folder{4}
\pages{1}{11}
\dating{[à partir de 1958]}
\watermark{Édition de démonstration\\interprétation personnelle de l'œuvre}
\foldertitle{Torsion en cohomologie cristalline : notes manuscrites (s.d.), copies de tapuscrit annoté (s.d.) — lecture modernisée du dossier entier}

\begin{document}

\section*{Résumé}

\begin{resume}
Une variété algébrique définie sur un corps fini a un nombre fini de points,
et ce nombre se lit sur sa cohomologie : il est une somme alternée de
puissances de nombres algébriques, les \emph{valeurs propres de Frobenius}.
Ces feuillets posent une question fine sur ces nombres. Si l'on les réduit
modulo le nombre premier $p$ qui est la caractéristique du corps, certains
s'annulent et d'autres survivent ; ceux qui survivent --- les « unités
$p$-adiques » --- sont exactement ce que la cohomologie voit encore quand on
ne garde que l'arithmétique modulo $p$. Combien y en a-t-il ?

On peut penser à un objet dont on ne regarde que l'ombre : la cohomologie
cristalline est l'objet entier, avec ses nombres ; la cohomologie
« cohérente » des fonctions sur la variété, que l'on calcule avec de
l'algèbre linéaire ordinaire en caractéristique $p$ (c'est la matrice de
Hasse-Witt), en est l'ombre portée. Le texte cherche à dire exactement
combien de l'objet passe dans l'ombre. La réponse est une suite exacte, le
« théorème » 8.5 : le nombre de valeurs propres unités est au moins le rang
de la partie de l'ombre qui « survit » à l'itération du Frobenius, et la
différence est mesurée par la \emph{torsion} de la cohomologie cristalline,
c'est-à-dire par ce que les nombres entiers $p$-adiques voient et que les
nombres rationnels ne voient pas.

La variété étudiée n'est pas quelconque : c'est une section $Y$ d'une
variété fixée $X$ par une hypersurface, et l'on ne regarde que la partie de
la cohomologie de $Y$ qui ne vient pas de $X$, la cohomologie
« évanescente ». Quand $Y$ varie, la réponse s'organise en une famille ;
la partie qui vient de la torsion de $X$ ne bouge pas, et celle qui bouge
est gouvernée par les matrices de Hasse-Witt des $Y$. Le texte finit par une
question : le groupe qui décrit comment cette partie mobile tourne sur
elle-même quand $Y$ parcourt la famille est-il aussi gros que possible ? Si
oui, on pourrait choisir $Y$ de façon que son nombre de points soit congru
modulo $p$ à n'importe quel résidu donné d'avance --- et le texte note que
personne ne savait le faire, pas même pour les courbes planes.

Le dossier montre aussi une manière de travailler. Le texte dactylographié
est relu à la main, et la relecture change la démonstration : là où le texte
passait par la dualité de Poincaré, une feuille ajoutée calcule directement,
et ce calcul rend fausses deux affirmations et un exemple que le texte
avançait, qui sont barrés. La marge le dit en une phrase : il est malcommode
d'utiliser ici la dualité. Le titre lui-même prévient que tout cela est
« heuristique » : plusieurs ingrédients de la cohomologie cristalline sont
admis, dont certains ont été démontrés depuis.

\keywords{crystalline cohomology, p-torsion, unit-root eigenvalues, vanishing cohomology, Hasse–Witt matrix, mod p monodromy}
\end{resume}

\section*{Le fil du dossier, et les conventions}

Le dossier tient en deux pièces : les pages 2 à 11, copies d'un tapuscrit
qui est le §8 d'un texte plus long --- il renvoie à des n°~6 et~7 qui ne sont
pas ici ---, repris au crayon et à l'encre ; et la page 1, une feuille
manuscrite marquée « 8.4.4 », qui est la « feuille jointe » à laquelle renvoie
la marge de la page 5. On lit le tapuscrit dans l'ordre, et la page 1 à la
place que la révision lui donne, après le lemme 8.4.3.\footnote{La datation
« à partir de 1958 » est celle de l'inventaire et on la garde. Les mentions de
Berthelot (page 2) et de Katz (pages 6 et 11) placent sans doute la rédaction
plus tard, dans la seconde moitié des années 1960 ; c'est une inférence de
notre part, non une datation.}

\begin{enumerate}
\item[1.] \emph{Le cadre} (page 2) : coefficients universels entre cohomologie
cristalline et cohomologie de De Rham ; la suite du serpent pour la
restriction de $X$ à $Y$.
\item[2.] \emph{Une voie abandonnée} (pages 2 et 3) : le n°~8.2, barré, qui
calculait les termes de torsion par dualité de Poincaré.
\item[3.] \emph{Lefschetz} (pages 3 et 4) : la cohomologie évanescente est
sans torsion.
\item[4.] \emph{Parties semi-simples} (pages 4 et 5) : le foncteur exact
$V \mapsto V_{\mathrm{ss}}$, et le lemme 8.4.3.
\item[5.] \emph{La feuille 8.4.4} (page 1, et la marge de la page 2) :
nilpotence de $F$ sur la torsion au-dessus de la dimension.
\item[6.] \emph{Le « théorème » 8.5} (pages 5 à 7) : la suite exacte (8.5.1)
et l'inégalité (8.5.6).
\item[7.] \emph{Remarques} (pages 7 à 10) : injectivité de $v$, variation
avec $Y$, critère de bijectivité de $\varphi$, l'exemple de Serre.
\item[8.] \emph{Monodromie} (page 11) : la question finale.
\end{enumerate}

Il y a donc deux couches, et elles ne disent pas la même chose. Le tapuscrit
calcule la torsion par dualité, et en tire des conditions sur
$\mathrm{Tors}^{n-1}$ et $\mathrm{Tors}^{n+1}$ ; la révision calcule
directement, trouve que $\varphi$ est toujours surjectif, et barre ce qui
disait le contraire. La lecture suit la révision, et dit en note ce que le
tapuscrit avait.

\emph{Conventions, valables pour tout le dossier.} $k$ est un corps parfait de
caractéristique $p > 0$, $W = W(k)$ l'anneau des vecteurs de Witt, $\sigma$
son Frobenius. $X \subset \mathbf{P}^{r}_{k}$ est lisse, projective, connexe,
de dimension $n$ ; $H^{i}(X)$ désigne sa cohomologie cristalline, un
$W$-module de type fini, muni d'un Frobenius $F$ $\sigma$-linéaire ;
$\mathrm{Tors}^{i}(X)$ est son sous-module de torsion,
$H^{i}_{\mathrm{lib}}(X) = H^{i}(X)/\mathrm{Tors}^{i}(X)$, et
${}_{p}\mathrm{Tors}^{i}(X) = \{x \in H^{i}(X) : px = 0\}
= \mathrm{Tor}^{W}_{1}(H^{i}(X), k)$. $Y$ est une section lisse de $X$ par une
hypersurface de degré $d \geq 1$. Pour un $k$-vectoriel de dimension finie
muni d'un endomorphisme $p$-linéaire $F$, l'indice $\mathrm{ss}$ désigne la
partie semi-simple pour $F$ ; l'exposant $\mathrm{ss}$, celle pour
l'opérateur transposé $V$, $p^{-1}$-linéaire. Ces deux conventions sont
celles de la page 1.

\emph{Ce que le dossier annonce et n'établit pas} : l'injectivité de $v$ en
général (8.5.7) ; la non-nullité générique de
$E(Y, \mathcal{O}_{Y})_{\mathrm{ss}}$ pour $d$ grand (8.5.8, renvoyée au
n°~7) ; la « conjecture du n°~6 » sur le niveau de la cohomologie
évanescente ; la réponse à la question de monodromie de 8.6, et la
congruence sur $\mathrm{card}\, Y(k')$ qui en découlerait. Le théorème de
Lefschetz faible, torsion comprise, est admis en cohomologie cristalline
d'un bout à l'autre.

\pagerange{2}{2}
\section*{Le cadre : coefficients universels et suite du serpent (page 2)}

\textbf{8.1.} La cohomologie cristalline et la cohomologie de De Rham sont
liées par l'isomorphisme, dans la catégorie dérivée des $k$-modules,
\[
(8.1.1) \qquad R\Gamma_{\mathrm{DR}}(X/k) \simeq R\Gamma_{\mathrm{cris}}(X/W)
\otimes^{L}_{W} k ,
\]
d'où, pour tout $i$, une suite exacte fonctorielle en $X$ :
\[
(8.1.2) \qquad 0 \to H^{i}(X) \otimes_{W} k \to H^{i}_{\mathrm{DR}}(X)
\to {}_{p}\mathrm{Tors}^{i+1}(X) \to 0 .
\]
\footnote{La page écrit « si on dispose d'une bonne théorie de la cohomologie
cristalline », et note que « ceci marche dès maintenant pour $p \neq 2$, grâce
à Berthelot » (un « Sauf erreur » est biffé devant). L'isomorphisme (8.1.1)
est démontré depuis pour tout $p$, $X$ propre et lisse sur $k$ parfait
(Berthelot, 1974 ; Berthelot et Ogus, \emph{Notes on crystalline cohomology},
1978). Le « $L$ » de (8.1.1) est porté à la main.}
Le Frobenius de $H^{i}(X)$ induit sur $H^{i}(X) \otimes_{W} k$ un endomorphisme
$p$-linéaire, et sur $H^{i}_{\mathrm{DR}}(X)$ c'est l'endomorphisme $F_{X}^{*}$
induit par le Frobenius absolu de $X$ ; (8.1.2) leur est compatible.

\emph{La restriction à $Y$.} Pour $i = n - 1$, la restriction de $X$ à $Y$
envoie la suite (8.1.2) de $X$ dans celle de $Y$, par un triple
$u = (u_{0} \otimes k, u_{1}, u_{2})$, où $u_{0} : H^{n-1}(X) \to H^{n-1}(Y)$,
$u_{1} : H^{n-1}_{\mathrm{DR}}(X) \to H^{n-1}_{\mathrm{DR}}(Y)$, et
\[
(8.1.4) \qquad u_{2} : {}_{p}\mathrm{Tors}^{n}(X) \to {}_{p}\mathrm{Tors}^{n}(Y) .
\]
Posons
\[
E(Y) = \mathrm{Coker}\bigl(H^{n-1}(X) \xrightarrow{u_{0}} H^{n-1}(Y)\bigr),
\qquad
E_{\mathrm{DR}}(Y) = \mathrm{Coker}\, u_{1} ,
\]
la cohomologie \emph{évanescente}, cristalline et de De Rham. Le produit
tensoriel étant exact à droite, $\mathrm{Coker}(u_{0} \otimes k) = E(Y)
\otimes_{W} k$, et le lemme du serpent donne la suite exacte
\[
(8.1.3) \qquad 0 \to \mathrm{Ker}(u_{0} \otimes k) \to \mathrm{Ker}\, u_{1}
\to \mathrm{Ker}\, u_{2} \to E(Y) \otimes_{W} k \to E_{\mathrm{DR}}(Y)
\to \mathrm{Coker}\, u_{2} \to 0 .
\]
\footnote{La page écrit « $\mathrm{Ker}\, u_{0}$ » pour le premier terme ;
c'est le noyau de $u_{0} \otimes k$ qui intervient. Le triplet $u$, la mise
entre parenthèses des trois premiers termes et « $= \mathrm{Coker}\, u_{1}$ »
sont de sa main. Le texte dit « section hyperplane », un « de degré assez
élevé » biffé ; les hypersurfaces de degré $d$ s'y ramènent par le plongement
de Veronese, comme le dira 8.5.}

\pagerange{2}{3}
\subsection*{La voie de la dualité, abandonnée (pages 2 et 3)}

Le n°~8.2 du tapuscrit, entièrement barré, interprétait les termes extrêmes
par dualité de Poincaré. Comme $\mathrm{Tor}^{W}_{1}(-, k) \simeq
\mathrm{Hom}_{W}(k, -)$, et que la dualité entre torsions échange
$\mathrm{Tors}^{i}$ et $\mathrm{Tors}^{2m+1-i}$ sur une variété de dimension
$m$, le but de $u_{2}$ est dual de $\mathrm{Tors}^{n-1}(Y) \otimes k$, sa
source de $\mathrm{Tors}^{n+1}(X) \otimes k$, et $u_{2}$ est le transposé de
l'application de Gysin
\[
(8.2.1) \qquad u'_{2} : \mathrm{Tors}^{n-1}(Y) \otimes_{W} k \to
\mathrm{Tors}^{n+1}(X) \otimes_{W} k ;
\]
$\mathrm{Ker}\, u_{2}$ et $\mathrm{Coker}\, u_{2}$ sont alors les duaux de
$\mathrm{Coker}\, u'_{2}$ et $\mathrm{Ker}\, u'_{2}$.\footnote{Ces dualités sont
correctes au niveau des modules, sous la dualité de Poincaré cristalline
admise. Une note de biais, en regard, dit qu'« il est \emph{malcommode}
d'utiliser ici la dualité de Poincaré » (le mot souligné est d'une lecture
incertaine). Ce qui est malcommode, à notre sens, est que la dualité transpose
$F$ : la partie $F$-semi-simple d'un dual est duale de la partie semi-simple
pour l'opérateur transposé, non pour $F$. C'est pourquoi la page 1 introduit
deux notations, indice et exposant. Les termes de torsion se calculent
directement (page 1, et plus bas), et la suite du texte n'a plus besoin de
8.2.}

\pagerange{3}{4}
\section*{Lefschetz : la cohomologie évanescente est sans torsion (pages 3 et 4)}

\textbf{8.3.} En cohomologie $\ell$-adique ($\ell \neq p$), le théorème de
Lefschetz faible à coefficients $\mathbf{Z}/\ell^{\nu}\mathbf{Z}$ dit que
$u_{0}$ est injectif et que son conoyau $E(Y)$ est \emph{sans
torsion}.\footnote{Le tapuscrit invoquait la situation complexe et le livre
de Lefschetz ; la main remplace par la cohomologie $\ell$-adique et le
théorème « faible ». L'énoncé avec torsion est juste : si $U = X - Y$, affine
de dimension $n$, on a $H^{i}_{c}(U) = 0$ pour $i < n$ et
$H^{n}_{c}(U, \mathbf{Z}_{\ell})$ sans torsion (sa torsion est duale de celle
de $H^{n+1}(U)$, nulle par dimension cohomologique), et $E(Y) \subset
H^{n}_{c}(U)$.} On \emph{admet} qu'il en va de même en cohomologie
cristalline.\footnote{C'est l'hypothèse la plus forte du texte. Nous ne
connaissons pas avec certitude l'état actuel de ce théorème de Lefschetz
entier, torsion comprise, pour la cohomologie cristalline, et ne
l'affirmons pas.}
Il en résulte :
\begin{enumerate}
\item[(i)] $u_{0}$ induit un isomorphisme $\mathrm{Tors}^{n-1}(X) \simeq
\mathrm{Tors}^{n-1}(Y)$ (8.3.1), puisque la torsion de $H^{n-1}(Y)$ s'envoie
sur zéro dans $E(Y)$ ;
\item[(ii)] $E(Y)$ est un $W$-module libre, de sorte que
$\mathrm{Tor}^{W}_{1}(E(Y), k) = 0$ et que, $u_{0}$ étant injectif,
\[
(8.3.3) \qquad \mathrm{Ker}(u_{0} \otimes k) = 0 ;
\]
\item[(iii)] les valeurs propres d'un endomorphisme de $E(Y)$ --- un
Frobenius --- se réduisent modulo $p$ en celles de l'endomorphisme induit sur
$E(Y) \otimes k$.
\end{enumerate}
\footnote{Le tapuscrit tire aussi de (8.3.1) que l'application (8.2.1) devient
le cup-produit par la classe de polarisation,
$u''_{2} : \mathrm{Tors}^{n-1}(X) \otimes k \to \mathrm{Tors}^{n+1}(X) \otimes
k$ (8.3.2), qui ne dépend plus de $Y$. Avec l'abandon de 8.2, ce
cup-produit ne sert plus. Une insertion à la main au-dessus d'« aussi du »
biffé ne se lit pas.}

\pagerange{4}{5}
\section*{Parties semi-simples, et le lemme 8.4.3 (pages 4 et 5)}

\textbf{8.4.} Soit $V$ un $k$-vectoriel de dimension finie muni d'un
endomorphisme $p$-linéaire $F$. Comme $k$ est parfait, les images
$F^{m}(V)$ sont des sous-espaces, et l'on a la décomposition
\[
V = V_{\mathrm{ss}} \oplus V_{\mathrm{nil}}, \qquad
V_{\mathrm{ss}} = \bigcap_{m} F^{m}(V), \quad
V_{\mathrm{nil}} = \bigcup_{m} \mathrm{Ker}\, F^{m} ,
\]
où $F$ est bijectif sur $V_{\mathrm{ss}}$ et nilpotent sur $V_{\mathrm{nil}}$.
$V_{\mathrm{ss}}$ est le plus grand sous-espace stable sur lequel $F$ est
surjectif ; le foncteur $V \mapsto V_{\mathrm{ss}}$ est
\emph{exact}.\footnote{C'est la décomposition de Fitting, pour un
endomorphisme semi-linéaire ; le nom est nôtre. Le tapuscrit la décrit en
prose, et souligne à la main « nilpotent ».} Le même énoncé vaut, avec
$\bigcap F^{m}(T)$, pour un $W$-module de longueur finie $T$ muni d'un
endomorphisme $\sigma$-linéaire ; on s'en servira pour la torsion.

Appliqué à (8.1.3), compte tenu de (8.3.3), le foncteur donne la suite
exacte
\[
(8.4.1) \qquad 0 \to (\mathrm{Ker}\, u_{1})_{\mathrm{ss}} \to
(\mathrm{Ker}\, u_{2})_{\mathrm{ss}} \to (E(Y) \otimes_{W} k)_{\mathrm{ss}}
\to E_{\mathrm{DR}}(Y)_{\mathrm{ss}} \to (\mathrm{Coker}\, u_{2})_{\mathrm{ss}}
\to 0 .
\]

\textbf{8.4.2.} Un vectoriel à opérateur $p$-linéaire bijectif est la même
chose qu'une représentation continue de $G = \mathrm{Gal}(\overline{k}/k)$ sur
un $\mathbf{F}_{p}$-vectoriel de dimension finie : à $V$ on associe
$(V \otimes_{k} \overline{k})^{F = 1}$, de même dimension, et l'on retrouve
$V$ par extension des scalaires. En famille, sur une base $S$, les Modules
localement libres à Frobenius bijectif correspondent de même aux systèmes
locaux étales de $\mathbf{F}_{p}$-vectoriels.\footnote{Le premier énoncé
repose sur le théorème de Lang ($F - 1$ est surjectif sur
$V \otimes \overline{k}$) ; la version en famille est classique, on la trouve
par exemple chez Katz (1973) pour les cristaux unités. Le tapuscrit ajoute une
autodualité de cette catégorie, valable dans les deux langages.}

\textbf{Lemme 8.4.3.} Soit $X$ un schéma propre sur $k$. Pour tout $i$,
l'application d'augmentation $H^{i}_{\mathrm{DR}}(X) \to H^{i}(X,
\mathcal{O}_{X})$ induit un isomorphisme sur les parties semi-simples pour
le Frobenius.\footnote{Le numéro tapé « 8.4.2 » est corrigé à la main.
L'énoncé dit « un corps $k$ de car.\ $p > 0$ » ; on garde $k$ parfait, comme
dans tout le texte. Pour $k$ quelconque, il suffit de remplacer les images
$F^{m}(V)$ par les sous-espaces qu'elles engendrent.}

\emph{Démonstration.} L'augmentation provient du quotient
$\Omega^{\bullet}_{X/k} \to \mathcal{O}_{X}$ du complexe de De Rham par son
sous-complexe $\sigma_{\geq 1}\Omega^{\bullet}$. Comme $F_{X}^{*}(dx) =
d(x^{p}) = 0$, le Frobenius est nul sur ce sous-complexe, donc sur sa
cohomologie, dont la partie semi-simple est nulle ; l'exactitude du foncteur,
appliquée à la suite longue, conclut.\footnote{La page passe par la suite
spectrale de Hodge $E_{1}^{a,b} = H^{b}(X, \Omega^{a}_{X/k})$, sur laquelle
Frob est nul pour $a \neq 0$, de sorte que les
$(E_{1}^{0,b})_{\mathrm{ss}}$ sont des cycles universels, et s'arrête sur
« on gagne \ldots{} ». L'argument donné ici est le même, sans la suite
spectrale.}

En combinant le lemme avec (8.1.2), qui reste exacte sur les parties
semi-simples, on obtient pour tout $i$
\[
0 \to (H^{i}(X) \otimes_{W} k)_{\mathrm{ss}} \to H^{i}(X,
\mathcal{O}_{X})_{\mathrm{ss}} \to ({}_{p}\mathrm{Tors}^{i+1}(X))_{\mathrm{ss}}
\to 0 .
\]
C'est la suite que la marge de la page 2 écrit « grâce à 8.4.3 » ; la
feuille 8.4.4 en tire ce dont le théorème a besoin.

\pagerange{1}{2}
\section*{La feuille 8.4.4 : nilpotence au-dessus de la dimension (page 1)}

\textbf{8.4.4.} (a) \emph{Pour $i > n$, $F$ est nilpotent sur
$H^{i}(X) \otimes_{W} k$, sur $\mathrm{Tors}^{i}(X) \otimes_{W} k$ et sur
${}_{p}\mathrm{Tors}^{i}(X)$.}

En effet $H^{i}(X, \mathcal{O}_{X}) = 0$ pour $i > n$, donc, par la suite
précédente, $(H^{i}(X) \otimes k)_{\mathrm{ss}} = 0$. Comme $H^{i}_{\mathrm{lib}}$
est libre, $\mathrm{Tors}^{i} \otimes k \subset H^{i} \otimes k$, et $F$ y est
nilpotent. Enfin, si $F^{m}(T) \subset pT$ pour $T = \mathrm{Tors}^{i}(X)$,
alors, $F$ étant $\sigma$-linéaire, $F^{jm}(T) \subset p^{j}T$, donc $F$ est
nilpotent sur $T$, et a fortiori sur ${}_{p}T$.\footnote{La page énonce (a)
sans démonstration ; l'argument est nôtre. Il donne en particulier
$({}_{p}\mathrm{Tors}^{n+1}(X))_{\mathrm{ss}} = 0$, qui ne découle pas
de la seule suite de coefficients universels. « $\mathrm{Tor}$ » et l'indice
$W$ de la ligne (a) sont d'une lecture incertaine ; on écrit
${}_{p}\mathrm{Tors}^{i}$, qui leur est isomorphe.}

Il en résulte, avec la suite précédente pour $i = n$, que
$H^{n}(X, \mathcal{O}_{X})_{\mathrm{ss}} \simeq (H^{n}(X) \otimes_{W}
k)_{\mathrm{ss}}$, d'où
\[
(*) \qquad 0 \to (\mathrm{Tors}^{n}(X) \otimes k)_{\mathrm{ss}}
\to H^{n}(X, \mathcal{O}_{X})_{\mathrm{ss}}
\to (H^{n}_{\mathrm{lib}}(X) \otimes k)_{\mathrm{ss}} \to 0 .
\]
\footnote{Les exposants et l'ordre des termes de $(*)$ sont lus sur une
première tentative effacée et ne sont pas assurés ; la forme donnée ici est
celle que l'argument établit. La marge de la page 2 note de même
« (nul si $i > n$) » sous $H^{i}(X, \mathcal{O}_{X})_{\mathrm{ss}}$ et
« (nul si $i \geq n$) », de lecture incertaine, au-dessus de
$({}_{p}\mathrm{Tors}^{i+1})_{\mathrm{ss}}$ : les deux sont justes.}

(b) Dualement, pour l'opérateur $V$ transposé de $F$ par la dualité de
Poincaré en cohomologie de De Rham --- sur $H^{\bullet}(X, \Omega^{n}_{X})$,
c'est l'opérateur de Cartier ---, $V$ est nilpotent sur
$H^{i}(X) \otimes k$ pour $i < n$, et sur $\mathrm{Tors}^{i}(X) \otimes k$ et
${}_{p}\mathrm{Tors}^{i}(X)$ pour $i \leq n$ ; et l'on a la suite transposée
de $(*)$ :
\[
(**) \qquad 0 \to (H^{n}_{\mathrm{lib}}(X) \otimes k)^{\mathrm{ss}}
\to H^{0}(X, \Omega^{n}_{X})^{\mathrm{ss}}
\to ({}_{p}\mathrm{Tors}^{n+1}(X))^{\mathrm{ss}} \to 0 .
\]
\footnote{La page dit seulement « sans doute $(*)$ et $(**)$ sont transposés
l'un de l'autre », suivi d'un mot illisible, et souligne deux fois
« $i \leq n$ ». La transposition est juste, terme à terme : $H^{0}(\Omega^{n})$
est dual de $H^{n}(\mathcal{O})$ par Serre, ${}_{p}\mathrm{Tors}^{n+1}$ de
$\mathrm{Tors}^{n} \otimes k$ par la dualité de Poincaré entre torsions, et
la dualité échange les deux parties semi-simples. Elle repose sur la dualité
de Poincaré cristalline, admise. La marge de la page 2 définit
$V_{X} = \sigma^{-1} F_{X*}$, qui est l'opérateur de type Cartier ; nous
prenons la transposition comme définition.} Plus généralement, pour tout
$i$, la marge de la page 2 écrit
\[
0 \to (H^{i}(X) \otimes_{W} k)^{\mathrm{ss}} \to
H^{i-n}(X, \Omega^{n}_{X})^{\mathrm{ss}}
\to ({}_{p}\mathrm{Tors}^{i+1}(X))^{\mathrm{ss}} \to 0 ,
\]
qui est (8.1.2) prise en parties $V$-semi-simples, compte tenu de ce que
$H^{i-n}(X, \Omega^{n}_{X})^{\mathrm{ss}} \simeq
H^{i}_{\mathrm{DR}}(X)^{\mathrm{ss}}$, transposé du lemme
8.4.3.\footnote{La page porte pour le dernier terme l'exposant $i - n + 1$, de
lecture incertaine ; c'est $i + 1$ qui est juste, comme le montre $(**)$ pour
$i = n$. L'exposant $i - n$ du terme médian est juste.}

\emph{Notation} : l'indice $\mathrm{ss}$ pour la partie semi-simple
relativement à $F$, l'exposant $\mathrm{ss}$ relativement à $V$.

\emph{Application aux termes de (8.4.1).} La variété $Y$ est de dimension
$n - 1 < n$ ; (a) appliqué à $Y$ rend $F$ nilpotent sur
${}_{p}\mathrm{Tors}^{n}(Y)$. Donc
\[
(\mathrm{Ker}\, u_{2})_{\mathrm{ss}} = ({}_{p}\mathrm{Tors}^{n}(X))_{\mathrm{ss}},
\qquad (\mathrm{Coker}\, u_{2})_{\mathrm{ss}} = 0 .
\]
C'est ce qu'écrit la marge de la page 4, « en vertu de
8.4.4 ».\footnote{Les deux égalités sont écrites au-dessus du terme et
encerclées en marge. À la page 5, le long du lemme, la marge renvoie à « Cor.\
8.4.4 \emph{ter} (utilisé dans 8.1.\,?) les résultats de la feuille jointe
(marquée 8.4.4) » ; la numérotation « ter » et le renvoi sont incertains.}
Et par le lemme 8.4.3, $(\mathrm{Ker}\, u_{1})_{\mathrm{ss}}$ et
$E_{\mathrm{DR}}(Y)_{\mathrm{ss}}$ sont le noyau et le conoyau de
\[
(8.5.2) \qquad v : H^{n-1}(X, \mathcal{O}_{X})_{\mathrm{ss}} \to
H^{n-1}(Y, \mathcal{O}_{Y})_{\mathrm{ss}} .
\]

\pagerange{5}{7}
\section*{Le « théorème » 8.5 (pages 5 à 7)}

Posons $E(Y, \mathcal{O}_{Y}) = \mathrm{Coker}\bigl(H^{n-1}(X,
\mathcal{O}_{X}) \to H^{n-1}(Y, \mathcal{O}_{Y})\bigr)$, la cohomologie
évanescente \emph{cohérente} ; par exactitude,
$E(Y, \mathcal{O}_{Y})_{\mathrm{ss}} = \mathrm{Coker}\, v$.

\textbf{« Théorème » 8.5.} \emph{Admettons, pour la cohomologie cristalline,
la formule des coefficients universels et le théorème de Lefschetz faible
torsion comprise. Soit $X \subset \mathbf{P}^{r}_{k}$ lisse projective connexe
de dimension $n$ sur un corps parfait de caractéristique $p > 0$, et $Y$ une
section de $X$ par une hypersurface de degré $d \geq 1$, lisse de dimension
$n - 1$. On a une suite exacte}
\[
(8.5.1) \qquad 0 \to \mathrm{Ker}\, v \xrightarrow{w}
({}_{p}\mathrm{Tors}^{n}(X))_{\mathrm{ss}} \to
(E(Y) \otimes_{W} k)_{\mathrm{ss}} \xrightarrow{\varphi}
E(Y, \mathcal{O}_{Y})_{\mathrm{ss}} \to 0 ,
\]
\emph{où $w$ est la restriction à $\mathrm{Ker}\, v \subset
H^{n-1}_{\mathrm{DR}}(X)_{\mathrm{ss}}$ de la surjection
$H^{n-1}_{\mathrm{DR}}(X)_{\mathrm{ss}} \to
({}_{p}\mathrm{Tors}^{n}(X))_{\mathrm{ss}}$ de (8.1.2).}

C'est (8.4.1), réécrite avec les quatre identifications qui précèdent.
\footnote{L'énoncé tapé admet aussi la dualité de Poincaré ; l'argument
révisé ne s'en sert pas, la page 1 ayant remplacé le calcul par dualité.
« Sans négliger la torsion » remplace à la main « “faible”, y compris pour
les groupes de torsion », biffé ; « de degré $d \geq 1$ » est de sa main.
La suite (8.5.1) est entièrement reprise à la main : le second terme tapé est
noirci et remplacé par $({}_{p}\mathrm{Tors}^{n}(X))_{\mathrm{ss}}$ encadré,
la fin de la ligne tapée est noircie, et les étiquettes $w$, $\varphi$ sont
ajoutées. La description explicite de $w$ est nôtre ; le tapuscrit définissait
sous ce nom, en (8.5.3), le cup-produit par la classe de polarisation $d\xi$
sur la torsion, $(\mathrm{Tors}^{n-1}(X) \otimes k)_{\mathrm{ss}} \to
(\mathrm{Tors}^{n+1}(X) \otimes k)_{\mathrm{ss}}$ (étiquette lue
$u_{d\xi}$, incertaine), et ce passage est barré. En marge, au crayon très
pâle, « Katz fait observer \ldots{} sur la », le reste illisible.}

\emph{Interprétation.} $E(Y)$ étant libre, si $k$ est fini, le rang de
$(E(Y) \otimes k)_{\mathrm{ss}}$ est le nombre --- avec multiplicités --- des
valeurs propres du Frobenius relatif à $k$ opérant sur la cohomologie
cristalline évanescente de $Y$ \emph{qui sont des unités $p$-adiques} ;
réduites modulo $p$, elles sont les valeurs propres du Frobenius linéarisé
sur $(E(Y) \otimes k)_{\mathrm{ss}}$, ou, dans le langage de 8.4.2, de
l'inverse du Frobenius arithmétique de $G$ sur le
$\mathbf{F}_{p}$-vectoriel correspondant.\footnote{Le décompte repose sur la
décomposition de Fitting de $E(Y)$ lui-même, en partie où $F$ est bijectif et
partie où il est topologiquement nilpotent ; c'est là que l'absence de
torsion sert. Le texte ajoute que ces valeurs propres sont aussi celles de
la cohomologie $\ell$-adique évanescente, « sous réserve de l'algébricité »
des opérations sur les deux cohomologies : c'est établi depuis (Katz et
Messing, 1974, à l'aide des conjectures de Weil démontrées par Deligne). Le
« $E$ » de $E^{n-1}(Y)$ est frappé sur un « $H$ ».}

Comme $w$ est injectif, on a
\[
(8.5.5) \qquad \dim (E(Y) \otimes_{W} k)_{\mathrm{ss}} =
\dim E(Y, \mathcal{O}_{Y})_{\mathrm{ss}} +
\bigl(\dim ({}_{p}\mathrm{Tors}^{n}(X))_{\mathrm{ss}}
- \dim \mathrm{Ker}\, v\bigr) ,
\]
la parenthèse étant $\geq 0$, et a fortiori, par la surjectivité de
$\varphi$,
\[
(8.5.6) \qquad \dim (E(Y) \otimes_{W} k)_{\mathrm{ss}} \geq
\dim E(Y, \mathcal{O}_{Y})_{\mathrm{ss}} :
\]
\emph{le nombre de valeurs propres de Frobenius sur la cohomologie
évanescente qui sont des unités $p$-adiques est au moins le rang semi-simple
de la cohomologie évanescente cohérente}, c'est-à-dire le rang stable de la
matrice de Hasse-Witt de $Y$ en dimension critique, prise modulo celle de
$X$.\footnote{Le tapuscrit écrivait, dans (8.5.5) et (8.5.6), des termes
correctifs $- \dim \mathrm{Ker}\, w$ et $- t^{n-1}(X, p)$ (« nombre de
$p$-coefficients de torsion »), que la main noircit ou barre : la
surjectivité de $\varphi$, acquise par la page 1, les rend inutiles. Un
« $\mathrm{Coker}\, w$ » biffé est remplacé à la main.}

\pagerange{7}{9}
\section*{Injectivité de $v$, et variation avec $Y$ (pages 7 à 9)}

\textbf{8.5.7.} Le texte pense que $v$ est toujours injectif, ce qui
simplifierait (8.5.1) et (8.5.5) ; c'est vrai en tout cas
\begin{enumerate}
\item[(i)] si $d$ est assez grand pour que $H^{n-1}(X, \mathcal{O}_{X}(-d))
= 0$ : la suite $0 \to \mathcal{O}_{X}(-d) \to \mathcal{O}_{X} \to
\mathcal{O}_{Y} \to 0$ rend alors $H^{n-1}(X, \mathcal{O}_{X}) \to
H^{n-1}(Y, \mathcal{O}_{Y})$ injectif, et a fortiori $v$ ; pour $n \geq 2$,
cette annulation a lieu pour $d \gg 0$ par dualité de Serre ;
\item[(ii)] si $n = 2$.\footnote{« si $n = 2$, ou » est ajouté à la main, sans
raison donnée. Une justification : pour $n = 2$, par 8.4.2 et la suite
d'Artin-Schreier, $v$ correspond après extension à $\overline{k}$ à
$H^{1}(\overline{X}, \mathbf{Z}/p) \to H^{1}(\overline{Y}, \mathbf{Z}/p)$,
injectif par le théorème de Lefschetz pour le groupe fondamental ($Y$
ample et connexe, SGA 2). L'argument est nôtre.}
\end{enumerate}
\footnote{Un second alinéa, qui regroupait dans (8.5.5) les termes
$\dim \mathrm{Coker}\, w - \dim \mathrm{Ker}\, w$ de l'ancien $w$, est encadré
et barré.}

\textbf{8.5.8.} Prenons pour $Y$ une section par une hypersurface de degré
$d$ « assez générale », c'est-à-dire telle que le rang de
$E(Y, \mathcal{O}_{Y})_{\mathrm{ss}}$ soit maximal --- condition ouverte, ce
rang étant celui d'une puissance itérée de la matrice de Hasse-Witt. Alors,
au terme $- \dim \mathrm{Ker}\, v$ près, nul pour $d$ grand par 8.5.7, le
second membre de (8.5.5) \emph{ne dépend pas de $Y$}. Si $k$ est fini, il y a
donc un ouvert $U_{d}$ de l'espace projectif des hypersurfaces de degré $d$
sur lequel le nombre de valeurs propres unités de Frobenius sur la
cohomologie évanescente de $Y_{\overline{k}}$ est constant, pour les points de
$U_{d}$ à valeurs dans les extensions finies de $k$.

Le texte veut montrer que ce nombre est \emph{non nul} pour $d$ grand. Cela
impliquerait --- et serait sans doute équivalent à --- que le motif
évanescent $E(Y)$ est de \emph{niveau} maximal $n - 1$, pour tous ces $Y$, ou
pour $Y$ générique. Par (8.5.6), il suffit de savoir que la valeur générique
de $\dim E(Y, \mathcal{O}_{Y})_{\mathrm{ss}}$ est $> 0$ pour $d \to
\infty$.\footnote{Le « niveau » est celui de Grothendieck (1969) : l'écart
maximal $p - q$ des nombres de Hodge. Qu'une valeur propre unité force le
niveau maximal est une forme de la minoration « Newton au-dessus de Hodge »
(Mazur, puis Ogus, sous des hypothèses de non-torsion). La réciproque n'est
pas vraie $Y$ par $Y$ ; pour $Y$ générique, c'est une question d'ordinarité,
établie pour les intersections complètes générales dans l'espace projectif
(Illusie, 1990) --- ce qui règle le cas $X = \mathbf{P}^{n}$. Pour $X$
quelconque, nous ne savons pas, et ne l'affirmons pas. Deux lignes et demie à
la main, sous la ligne tapée, sont biffées : « on sait seulement (cf.\ n°~7)
que $\dim E(Y, \mathcal{O}_{Y})_{\mathrm{ss}} \to$ [illisible] pour $d$
grand », et que $\varphi \neq 0$ dès qu'une condition sur
$({}_{p}\mathrm{Tors}^{n+1}(X) \otimes k)_{\mathrm{ss}}$ est remplie ; la
marge renvoie : « or cela est établi au n°~7 ». Le n°~7 n'est pas dans le
dossier. Par 8.4.4 (a), cette partie semi-simple est d'ailleurs toujours
nulle.}

\textbf{8.5.9.} Soit $k$ à nouveau quelconque, et $U_{d}$ l'ouvert des
hypersurfaces de degré $d$ dont l'intersection $Y$ avec $X$ est lisse de
dimension $n - 1$ et pour lesquelles $E(Y, \mathcal{O}_{Y})_{\mathrm{ss}}$
est de rang maximal. Les suites (8.5.1), pour $Y$ parcourant $U_{d}$,
proviennent d'une suite exacte de Modules localement libres sur $U_{d}$ à
Frobenius bijectif, c'est-à-dire (8.4.2) d'une suite exacte de systèmes
locaux étales de $\mathbf{F}_{p}$-vectoriels sur $U_{d}$, ou de
représentations de $\pi_{1}(U_{d})$. Le noyau de $\varphi$, image de
$({}_{p}\mathrm{Tors}^{n}(X))_{\mathrm{ss}}$, est un système local
\emph{constant}, venu du corps de base.\footnote{La main biffe « et le
conoyau sont des » : $\varphi$ est surjectif. « Même pas nécessairement
parfait » est au tapuscrit ; les parties semi-simples se prennent alors comme
en note au lemme 8.4.3.}

Si $k$ est fini, quitte à l'étendre, le Frobenius relatif à $k$ a toutes
ses valeurs propres égales à $1$ sur ce noyau ;\footnote{La page dit qu'il
suffit qu'il en soit ainsi sur « $H^{n}(X, \mathcal{O}_{X})_{\mathrm{ss}}$ »,
l'exposant étant repris sur un chiffre noirci et d'une lecture incertaine. Ce
qui sert est $({}_{p}\mathrm{Tors}^{n}(X))_{\mathrm{ss}}$, quotient de
$H^{n-1}(X, \mathcal{O}_{X})_{\mathrm{ss}}$ ; c'est l'un ou l'autre qu'il faut
lire. « et sur les deux membres de (8.5.3) » est biffé.} la variation des
valeurs propres modulo $p$ du Frobenius sur $E(Y)$, quand $Y$ parcourt
$U_{d}$, est alors décrite par la seule représentation de $G =
\pi_{1}(U_{d})$ sur le $\mathbf{F}_{p}$-vectoriel des
$E(Y, \mathcal{O}_{Y})_{\mathrm{ss}}$ --- c'est-à-dire par les matrices de
Hasse-Witt des $Y$ en dimension critique ---, et par la répartition des
éléments de Frobenius dans $G$.

\pagerange{9}{10}
\section*{Quand $\varphi$ est-il bijectif ? L'exemple de Serre (pages 9 et 10)}

\textbf{8.5.10.} Par l'exactitude de (8.5.1), $\varphi$ est injectif, donc
bijectif, \emph{si et seulement si $w$ est surjectif}. Si $\mathrm{Ker}\, v =
0$, cela revient à $({}_{p}\mathrm{Tors}^{n}(X))_{\mathrm{ss}} = 0$ ; en tout
cas, il suffit que $X$ n'ait pas de $p$-torsion cristalline en degré $n$, ou,
ce qui revient au même par dualité, en degré $n + 1$.\footnote{La main biffe
partout les « resp. » qui faisaient aussi un énoncé de surjectivité, et
insère « donc bijectif » et « et, [illisible] si
$({}_{p}\mathrm{Tors}^{n}(X))_{\mathrm{ss}} = 0$ ». La page place la réserve
« $\mathrm{Ker}\, v = 0$ » sur le « seulement si » de l'équivalence avec la
surjectivité de $w$ ; cette équivalence-là est inconditionnelle, et la
réserve porte sur le critère par la torsion. L'exposant de la dernière
parenthèse, $n + 1$, est repris sur un chiffre noirci.}

Lorsque $X$ se relève, avec sa polarisation, en caractéristique nulle, cette
condition devrait pouvoir s'exprimer par la condition analogue sur la torsion
de la cohomologie $p$-adique de la variété relevée.\footnote{On sait
aujourd'hui que les torsions cristalline et étale $p$-adique se comparent
exactement pour les degrés petits devant $p$ (Fontaine et Messing ;
Caruso), et qu'en général seule une inégalité subsiste, la torsion étale
étant majorée par la torsion cristalline (Bhatt, Morrow et Scholze, 2018). La
suite, qui réécrivait l'ancien (8.5.3) en termes de cette variété, est barrée
de zigzags.}

\emph{Ce qui est barré, et pourquoi.} Le tapuscrit poursuivait : même pour
$n = 2$, $\varphi$ peut ne pas être surjectif ; puis il donnait (page 10)
l'équivalence de trois conditions --- $\varphi$ surjectif (resp.\ injectif)
pour tout $d$ et tout point de $U_{d}$ ; la même chose pour un seul $d$
multiple de $p$ avec $\mathrm{Ker}\, v = 0$ ; et
$(\mathrm{Tors}^{n-1}(X) \otimes k)_{\mathrm{ss}} = 0$ (resp.\
$(\mathrm{Tors}^{n+1}(X) \otimes k)_{\mathrm{ss}} = 0$). Tout cela est barré,
et à juste titre : par 8.4.4, $\varphi$ est toujours
surjectif.\footnote{La rature court de « on réécrit » au bas de la page 9
jusqu'à la condition c) de la page 10 ; qu'elle comprenne c) n'est pas sûr.
À notre sens, la condition de surjectivité de c) venait du calcul par
dualité, et son terme $\mathrm{Tors}^{n+1}$ y est toujours nul pour $F$ ; la
condition d'injectivité, elle, devient juste si l'on y lit la partie
semi-simple pour le transposé $V$, puisque
$({}_{p}\mathrm{Tors}^{n}(X))_{\mathrm{ss}}$ est duale de
$(\mathrm{Tors}^{n+1}(X) \otimes k)^{\mathrm{ss}}$. C'est exactement la
distinction que la page 1 introduit.}

\textbf{8.5.11.} Exemples où $\varphi$ n'est pas bijectif. Par (8.1.2) et le
lemme 8.4.3, pour tout $i$,
\[
\begin{aligned}
(8.5.11.1) \qquad \dim_{k} H^{i}(X, \mathcal{O}_{X})_{\mathrm{ss}} &=
\dim_{\mathbf{F}_{p}} H^{i}(\overline{X}, \mathbf{Z}/p\mathbf{Z}) \\
&= \dim (H^{i}_{\mathrm{lib}}(X) \otimes k)_{\mathrm{ss}}
+ \dim (\mathrm{Tors}^{i}(X) \otimes k)_{\mathrm{ss}} \\
&\qquad + \dim ({}_{p}\mathrm{Tors}^{i+1}(X))_{\mathrm{ss}} \\
&\geq \dim (H^{i}_{\mathrm{lib}}(X) \otimes k)_{\mathrm{ss}} ,
\end{aligned}
\]
la première égalité par la suite d'Artin-Schreier sur $\overline{X}$. Pour
$i = n - 1$ et $\mathrm{Ker}\, v = 0$, l'égalité a lieu si et seulement si
$\varphi$ est bijectif \emph{et} $(\mathrm{Tors}^{n-1}(X) \otimes
k)_{\mathrm{ss}} = 0$.\footnote{La page écrit l'inégalité sans la ligne
médiane et dit que « $\varphi$ bijectif signifie que l'inégalité précédente
est une égalité ». C'est vrai si l'on ne retranche pas la torsion au second
membre (avec $(H^{n-1}(X) \otimes k)_{\mathrm{ss}}$) ; avec
$H^{n-1}/\mathrm{Tors}^{n-1}$, comme écrit, l'égalité demande en outre
l'annulation de $(\mathrm{Tors}^{n-1} \otimes k)_{\mathrm{ss}}$. La
remarque « cela marche en toute dimension $i$ » est juste. Le début de
l'alinéa, « Il serait intéressant de donner des exemples où ces conditions
ne sont pas satisfaites \ldots{} », est biffé et remplacé par « $\varphi$
n'est pas bijectif ».}

\emph{L'exemple de Serre.} Serre, au symposium de Mexico, donne une variété
$X$ quotient d'une intersection complète par un groupe
$\mathbf{Z}/p\mathbf{Z}$ opérant sans point fixe, pour laquelle
$H^{1}(X) = 0$ en cohomologie cristalline, mais
$H^{1}(\overline{X}, \mathbf{Z}/p\mathbf{Z}) \simeq \mathbf{Z}/p\mathbf{Z}
\neq 0$.\footnote{La formule est coupée par le bord du feuillet après
« $\neq \mathbf{Z}/$ », que la main complète en « $p\mathbf{Z}$ ». Ce qu'a
l'exemple : l'intersection complète étant simplement connexe, le quotient a
pour groupe fondamental $\mathbf{Z}/p$, d'où $H^{1}(\overline{X},
\mathbf{Z}/p) \simeq \mathbf{Z}/p$ ; et $b_{1} = 0$, d'où $H^{1}(X) = 0$
($H^{1}$ cristallin étant toujours sans torsion). L'article est J.-P. Serre,
« Sur la topologie des variétés algébriques en caractéristique $p$ »,
Mexico, 1958.} En coupant par des hyperplans, on a le même exemple avec
$\dim X = 2$. Alors, par (8.1.2) pour $i = 1$, $H^{1}_{\mathrm{DR}}(X) \simeq
{}_{p}\mathrm{Tors}^{2}(X)$, et
\[
\dim ({}_{p}\mathrm{Tors}^{2}(X))_{\mathrm{ss}} =
\dim_{\mathbf{F}_{p}} H^{1}(\overline{X}, \mathbf{Z}/p\mathbf{Z}) = 1 ;
\]
comme $\mathrm{Ker}\, v = 0$ pour $n = 2$ (8.5.7), $w = 0$ et
$\mathrm{Ker}\, \varphi \simeq ({}_{p}\mathrm{Tors}^{2}(X))_{\mathrm{ss}} \neq
0$, quels que soient $d$ et $Y$ : \emph{$\varphi$ n'est pas
injectif}.\footnote{L'indice $p$, l'exposant $2$ repris sur un autre chiffre
et « $\mathrm{Ker}\, \varphi$ » sont de sa main ; au-dessus, il écrit
« $H^{1}(Y, \mathcal{O}_{Y})_{\mathrm{ss}}$ » avec un signe d'égalité
vertical dont le terme d'attache n'est pas assuré, et que nous ne
reprenons pas.}

\emph{Un exemple barré.} Le tapuscrit construisait ensuite, avec
$X' = X \times Z$ ($Z$ lisse projective de dimension $n - 2$, sans torsion,
$(H^{n-4}(Z) \otimes k)_{\mathrm{ss}} \neq 0$, et la formule de Künneth),
des exemples où $(\mathrm{Tors}^{n-1}(X') \otimes k)_{\mathrm{ss}} \neq 0$,
censés donner $\varphi$ non surjectif pour tout $n \geq 4$. Le paragraphe est
barré : $\varphi$ est toujours surjectif. Le calcul de torsion, lui, garde un
sens par (8.5.11.1) : il donne des exemples où l'inégalité est stricte pour
une autre raison que la non-bijectivité de $\varphi$.\footnote{Le « $4$ » est
repris à la main sur un « $3$ ». La dernière phrase est notre lecture de ce
qui reste valable dans l'exemple.}

\pagerange{11}{11}
\section*{La question de monodromie (page 11)}

\textbf{8.6.} La suite exacte de systèmes locaux sur $U_{d}$ issue de 8.5.9
montre que le système local des $(E(Y) \otimes_{W} k)_{\mathrm{ss}}$ a un
sous-système de nature triviale --- constant, venu du corps de base --- et
pour quotient celui des $E(Y, \mathcal{O}_{Y})_{\mathrm{ss}}$, dont on
s'attend qu'il donne une représentation tout à fait non triviale de $G =
\pi_{1}(U_{d})$ sur un $\mathbf{F}_{p}$-vectoriel $V$.\footnote{La main
corrige le tapuscrit, qui parlait d'un sous-système \emph{et} d'un quotient
triviaux et d'un composant commun « $\mathrm{Coim}\, \varphi =
\mathrm{Im}\, \varphi$ » : avec $\varphi$ surjectif, il n'y a plus de
quotient trivial. Le texte note que l'on peut tout lire au point générique
$s$ de l'espace des hypersurfaces, avec $\mathrm{Gal}(\overline{k(s)}/k(s))$
pour groupe, « mais on perd alors les frobenius ».}

\emph{Questions.} Pour $n \geq 2$, cette représentation est-elle
\emph{irréductible} ? L'image de $G$ dans $\mathrm{Aut}(V)$ est-elle même
$\mathrm{GL}(V)$ tout entier ? Dans ce dernier cas, dès que $\dim V \geq 2$
--- ce qui doit avoir lieu pour $d$ grand, $\dim V$ tendant vers l'infini
avec $d$ (8.5.8) ---, tout élément de $\mathbf{F}_{p}$ est la trace d'un
élément de $G$ opérant sur $V$, car l'application trace
$\mathrm{GL}_{m}(\mathbf{F}_{p}) \to \mathbf{F}_{p}$ est surjective pour
$m \geq 2$.\footnote{Les insertions qui transforment la question en
« irréductible, ou même » sont au crayon, « irréductible » souligné ; « Si
oui » est biffé. La justification par la trace est nôtre. Une longue note au
crayon en interligne commence par « (Ce n'est pas [illisible] lorsque
$X = \mathbf{P}^{r-1}_{k}$, cas qu'on peut traiter trivialement.) Lorsque la
condition [illisible] » ; une reprise en dessous ne se lit pas.}

\emph{Conséquence espérée.} Si $k$ est fini, la formule des traces modulo
$p$ (« formule de Woods Hole ») exprime $\mathrm{card}\, Y(k')$ modulo $p$
comme somme alternée des traces du Frobenius sur les
$H^{i}(Y, \mathcal{O}_{Y})$, donc sur leurs parties semi-simples ; la seule
contribution qui varie avec $Y$ est la trace sur $V$. Avec le théorème de
densité de Čebotarev, on en déduirait que, pour tout ouvert non vide $U$ de
l'espace des hypersurfaces de degré $d$ ($d$ comme ci-dessus), on peut
trouver un point $s$ de $U$ à valeurs dans une extension finie $k'$ de $k$
dont la section $Y$ a un nombre de points $\mathrm{card}\, Y(k')$ congru
modulo $p$ à un élément de $\mathbf{F}_{p}$ donné à
l'avance.\footnote{La formule est démontrée depuis pour $X$ propre sur un
corps fini (SGA~4$\frac{1}{2}$, « Fonctions $L$ modulo $\ell^{n}$ et modulo
$p$ » ; Katz, SGA 7, exposé XXII). Que seule la trace sur $V$ varie, pour $d$
grand, est notre explicitation. La déduction, telle que la page l'esquisse,
est heuristique : il faut aussi contrôler la monodromie arithmétique, et non
seulement géométrique.}

Le texte note qu'au moment d'écrire, et malgré l'aide de Katz, aucun
résultat de cette nature ne semblait prouvé, même pour $X = \mathbf{P}^{2}$,
c'est-à-dire pour les courbes planes $Y$ de grand degré.\footnote{Au-dessus
de $X = \mathbf{P}^{2}$, il écrit ce qui semble être « $2 = r$ ». Des
résultats de grande monodromie modulo $p$ ont été obtenus depuis pour
l'espace des modules des courbes (Ekedahl ; Achter et Pries) ; nous ne
savons pas s'ils couvrent les familles de sections hypersurfaces d'une $X$
donnée, et ne l'affirmons pas. Les deux dernières lignes et demie, qui
disaient qu'un tel résultat impliquerait la « conjecture du n°~6 sur le
niveau de la cohomologie évanescente » (cf.\ 8.5.8), sont encadrées et
barrées. Le n°~6 n'est pas dans le dossier.}

\end{document}
